%!TEX root = ../../thesis.tex

\section{Action Space}\label{sec:corpus-action-space}

% TODO: The composite action and its encoding --- the section where constraints C1,
% C2 and C5 of \cref{sec:mdp-constraints} actually bite.
%
% The factors, and why each is a decision rather than a heuristic:
%   - retain in {0..N} u {discard}: what to do with the PREVIOUS step's mutant,
%   - select in {0..N}: which corpus slot to mutate next,
%   - position in {0..L} and operator in the operator set.
%
% Argue the ordering: retain acts on the previous mutant because its coverage is
% already in the observation, so every eviction --- the only destructive decision --- is
% informed rather than blind. Note the pass-through case retain == select, which
% yields a hill-climbing lineage, and its cost (a lineage occupies a slot and
% overwrites its own parent).
%
% Encoding: pack the tuple into one integer of size (N+1) * N * L * |ops|. Give the
% configuration actually used and the resulting branching factor. Then discuss the
% alternative that was NOT taken and why:
%   - flat encoding keeps the policy a single softmax, so joint actions are
%     unrelated categories and nothing generalizes across factors,
%   - factored heads with sampled search (Sampled MuZero) would fix that but require
%     changing the policy model, the prior layout and the identity assertion in the
%     MCTS; deferred until the product outgrows enumeration,
%   - a subtlety worth stating once: the encoding *dimension* is additive (a
%     concatenation of one-hots) while the number of distinct joint actions --- the
%     policy width and the branching factor --- is the product. They are different
%     numbers about the same action.
%   - the flat encoding also represents the retain == select diagonal directly, which
%     independent factored heads cannot.
